\documentclass[11pt]{article}
\usepackage[margin=27mm]{geometry}
\usepackage{amsmath,amssymb,amsthm,booktabs,hyperref}
\newtheorem{lemma}{Lemma}
\newtheorem{theorem}[lemma]{Theorem}
\newcommand{\PSL}{\operatorname{PSL}}
\newcommand{\PGL}{\operatorname{PGL}}
\newcommand{\Aut}{\operatorname{Aut}}
\newcommand{\N}{\operatorname{N}}
\newcommand{\Cen}{\operatorname{C}}
\title{Pairs of symmetric parabolics and their full graph groups}
\author{Human Compression Project}\date{}
\begin{document}\maketitle
\begin{abstract}
For prime $p\equiv\pm1\pmod8$, the normalizer of a dihedral subgroup
of order eight determines the two $S_4$ parabolics above it. According
to one parity, these parabolics belong to the same conjugacy class or
to different classes. This distinction gives the full graph groups
$\PSL(2,p)\times C_2$ and $\PGL(2,p)$, respectively. The proof uses
a concrete type interchange and Tutte's cubic stabilizer bound; it
requires no graph automorphism search. The source's conjugacy counts
are recovered, while four of its nine prime-field full-group
identifications require correction.
\end{abstract}

\section{The normalizer determines the parabolics}
Let $p\ge7$ be prime with $p\equiv\pm1\pmod8$. Choose
$\epsilon\in\{1,-1\}$ so that $p\equiv\epsilon\pmod8$, and put
\[
 G=\PSL(2,p),\quad E=\PGL(2,p),\quad
 n=\frac{p-\epsilon}{2},\quad d=\frac{p-\epsilon}{8}.
\]
Thus $n=4d$. We retain the classical subgroup facts used by the source:
the maximal subgroups $S_4$ of $G$ form two conjugacy classes, fused
by $E$; their Sylow two-subgroups are $D_8$; cyclic groups of order
four lie in a unique torus of order $n$ in $G$ and $2n$ in $E$;
and those torus normalizers are dihedral. Also $G$ is simple and
$\Aut(G)=E$. The cyclic torus containing an element of order four
is its projective centralizer, so its uniqueness here is literal,
not just an abstract order identification.

\begin{lemma}\label{normalizer}
For $B\le G$ isomorphic to $D_8$,
\[
 \N_E(B)\cong D_{16},\qquad
 \N_G(B)=
 \begin{cases}D_{16},&d\text{ even},\\ B,&d\text{ odd}.
 \end{cases}
\]
The subgroups $B$ form one $E$-class. This class splits into two
$G$-classes when $d$ is even and remains one class when $d$ is odd.
\end{lemma}
\begin{proof}
The cyclic subgroup $C_4$ is characteristic in $B$, so every normalizer
of $B$ normalizes its unique torus. Write the torus normalizer in $E$
as
\[
 \langle R,S\mid R^{2n}=S^2=1,\ SRS=R^{-1}\rangle,
 \qquad \N_G(C_4)=\langle R^2,S\rangle,
\]
choosing a reflection $S\in G$. After conjugacy in this normalizer
we may take $B=\langle R^{n/2},S\rangle$. Indeed, its reflections
in $G$ have even exponents $R^{2j}S$, and conjugation by $R^j$
is transitive on these possibilities modulo the rotation subgroup of
$B$. Together with conjugacy of the tori this proves that the $D_8$
subgroups in $G$ form one $E$-class.

A rotation $R^k$ normalizes $B$ exactly when $(n/2)\mid2k$, or
$d\mid k$. Its normalizing rotations are therefore
$\langle R^d\rangle$, of order eight. The full normalizer is
$\langle R^d,S\rangle\cong D_{16}$. Since $R^k\in G$ exactly
when $k$ is even, this normalizer lies in $G$ when $d$ is even.
When $d$ is odd, its intersection with $G$ is
$\langle R^{2d},S\rangle=B$. For a subgroup of index two, an
$E$-conjugacy class splits into two $G$-classes precisely when its
normalizer is contained in $G$: this follows directly by comparing
the $G$-orbit and $E$-orbit lengths. This proves the final assertion.
\end{proof}

\begin{lemma}\label{overgroups}
Each $B\cong D_8$ is contained in exactly two subgroups $M_0,M_1$
isomorphic to $S_4$. Their intersection is $B$. If $d$ is even,
they belong to the same $G$-class and are interchanged by
$\N_G(B)/B$. If $d$ is odd, they belong to different $G$-classes
and are interchanged by $\N_E(B)/B$.
In either case there is an element $t$ interchanging them by conjugation
with $t^2\in B$; one may take $t\in G$ exactly in the even case.
\end{lemma}
\begin{proof}
Every $S_4$ contains one conjugacy class of three subgroups $D_8$,
and $\N_{S_4}(D_8)=D_8$. Fix an $S_4$ class containing the given
$B$. Its members containing $B$ form a transitive $\N_G(B)$-set:
adjust a transporter inside $S_4$ so that it also preserves $B$.
The stabilizer of one member is
$\N_G(B)\cap\N_G(S_4)=B$, since a maximal $S_4$ is
self-normalizing in the simple group $G$. Thus the number of such
members in this class is $[\N_G(B):B]$.

For even $d$, the two $B$-classes are exchanged by the outer coset,
as are the two $S_4$ classes. Each $S_4$ class contains only one
$B$-class; the two $S_4$ classes consequently contain different
$B$-classes. A fixed $B$ therefore has two overgroups in one class
and none in the other. Their interchange is the regular action of
$\N_G(B)/B\cong C_2$.

For odd $d$, there is just one $B$-class. Each of the two $S_4$
classes supplies one overgroup. Every element of $\N_E(B)\setminus G$
interchanges these overgroups, because it exchanges their $G$-classes.
In both cases the relevant normalizer is $D_{16}$ with $B$ of index
two, so the square of an interchanging element belongs to $B$.
Finally $B$ is maximal in $S_4$; distinct equal-order overgroups
containing it therefore meet exactly in $B$.
\end{proof}

These pairs give connected cubic bipartite coset graphs. The local
action of $S_4$ on its three cosets of $D_8$ is $S_3$, and distinct
maximal parabolics generate $G$. Thus all the rank two geometry and
local two-transitivity conditions hold.

\begin{theorem}\label{counts}
Under the source's unordered convention, the configurations with
parabolics $S_4,S_4$ and intersection $D_8$ have two conjugacy classes
when $d$ is even and one when $d$ is odd. In both cases there is one
class under $\Aut(G)$.
\end{theorem}
\begin{proof}
For each $B$ the unordered pair of overgroups is unique, by
Lemma~\ref{overgroups}. Its intersection recovers $B$. Conjugacy
and automorphism classes of pairs therefore correspond exactly to
the respective classes of $B$ counted in Lemma~\ref{normalizer}.
This proves the result, including the type interchange identification.
\end{proof}

\section{A literal type interchange}
Use left cosets $gM_i$ as vertices of type $i$, with incidence defined
by nonempty intersection. Left multiplication gives the specified
type-preserving action of $G$. Let $t$ be as in
Lemma~\ref{overgroups} and write $\alpha(g)=tgt^{-1}$.

\begin{lemma}\label{interchange}
The map
\[
 T(gM_i)=\alpha(g)M_{1-i}
\]
is a graph automorphism interchanging the two types. It satisfies
$T^2=L_{t^2}$, where $L_g$ denotes left multiplication by $g\in G$.
If $t\in G$, then $C=L_{t^{-1}}T$ is an involution centralizing
the whole action of $G$.
\end{lemma}
\begin{proof}
Conjugation by $t$ interchanges the two parabolics, so the displayed
formula is independent of the coset representative. It carries
intersecting cosets to intersecting cosets and is bijective, so it
is a graph automorphism. Since $t^2\in B\subseteq M_i$,
\[
 T^2(gM_i)=t^2gt^{-2}M_i=t^2gM_i=L_{t^2}(gM_i).
\]
When $t\in G$, the adjusted map has the simpler expression
$C(gM_i)=gt^{-1}M_{1-i}$. It commutes with every left multiplication,
and its square is $gM_i\mapsto gt^{-2}M_i=gM_i$. It is nonidentity
because it exchanges the two types.
\end{proof}

\section{The full automorphism group}
We retain Tutte's theorem for finite connected simple cubic
arc-transitive graphs: the full vertex stabilizer has order at most
$48$. The source and the separate subdivision proof state this
classical premise with its original references. Our graph is simple
by its coset-incidence definition and connected by maximality.
Lemma~\ref{interchange}, together with the local action of $G$,
gives arc transitivity of its full automorphism group.

\begin{theorem}\label{full}
The full automorphism group $A$ of this incidence graph is
\[
 A\cong
 \begin{cases}
 \PSL(2,p)\times C_2,& d\text{ even},\\
 \PGL(2,p),& d\text{ odd}.
 \end{cases}
\]
In both cases its exact arc-transitivity length is four.
\end{theorem}
\begin{proof}
There are $2|G|/24$ vertices and $A$ is vertex transitive. Tutte's
bound gives
\[
 [A:G]=\frac{2|A_x|}{24}\le4.
\]
The kernel of the action of $A$ on its cosets of $G$ is a normal
subgroup of $A$ contained in $G$, hence a normal subgroup of the
simple group $G$. If this kernel were trivial, $G$ would embed in
$S_4$, impossible since $|G|\ge168$. Therefore the kernel is $G$,
and $G\triangleleft A$.

An element of $\Cen_A(G)$ either preserves both $G$-orbits of
vertices or exchanges them. If it preserves both, it is the identity
on each: the centralizer of a transitive coset action $G/M_i$ is
$\N_G(M_i)/M_i=1$. This equality also follows directly by sending
the base vertex to another vertex fixed by its stabilizer. Thus the
centralizer has order at most two. If $M_0,M_1$ are in different
$G$-classes, it cannot exchange the orbits, because a $G$-equivariant
bijection of transitive $G$-sets requires conjugate point stabilizers.

For even $d$, Lemma~\ref{interchange} supplies precisely the
nontrivial centralizer element $C$. Conjugation embeds
$A/\langle C\rangle$ into $\Aut(G)=E$. Its image preserves the
single $G$-class containing both vertex stabilizers. Every outer
automorphism exchanges the two $S_4$ classes, so this image is
$\operatorname{Inn}(G)$. Since $G$ is centerless, it intersects
$\langle C\rangle$ trivially, and consequently
$A=G\times\langle C\rangle$.

For odd $d$, the centralizer is trivial and conjugation embeds $A$
in $E$. The constructed type interchange $T$ is outside the given
type-preserving $G$. Therefore $|A|\ge2|G|=|E|$, proving equality
with $E$. This proves both full-group assertions.

In either case $|A|=2|G|$, so $|A_x|=24=|G_x|$. The previously
proved local length four for the specified $S_4,S_4,D_8$ action
therefore remains the exact length for $A$: it has the same vertex
stabilizers. The type interchange makes this a global four-arc
transitive action. A fifth arc would require a vertex stabilizer
divisible by $3\cdot2^4=48$, which is impossible.
\end{proof}

\section{Comparison with the fixed source}
Theorem~\ref{counts} recovers Table 3, rows 7--8, including its
unordered convention. Theorem~\ref{full} applies to the nine prime
fields listed on printed page 386 of the fixed source, as follows.
\begin{center}
\begin{tabular}{rcll}
\toprule
$p$&$d$&Full graph group&Printed identification\\\midrule
17&2&$\PSL(2,17)\times C_2$&corrected\\
23&3&$\PGL(2,23)$&recovered\\
31&4&$\PSL(2,31)\times C_2$&corrected\\
41&5&$\PGL(2,41)$&recovered\\
47&6&$\PSL(2,47)\times C_2$&corrected\\
71&9&$\PGL(2,71)$&recovered\\
73&9&$\PGL(2,73)$&recovered\\
79&10&$\PSL(2,79)\times C_2$&corrected\\
89&11&$\PGL(2,89)$&recovered\\\bottomrule
\end{tabular}
\end{center}
The source identifies all nine full groups as $\PGL(2,p)$. In the
four corrected cases our explicit central involution already disproves
that identification: $\PGL(2,p)$ has trivial center, since its
centralizer of the normal simple subgroup $\PSL(2,p)$ is trivial.
The upper-bound argument proves the exact replacement, not just the
failure of the printed name. There is no second principal-refutation
credit: this is a surviving-claim correction in an original already
registered as principally refuted. The prime-square case $q=9$ is
separate and is not covered by this proof.

The geometry family and its source counts are attributed to the
original authors and the earlier work cited there. We retain the
classical subgroup structure, outer action and Tutte's bound. The
contribution here is the parity-controlled normalizer argument, the
explicit centralizing type interchange and the resulting corrected full
group theorem. The project root supplied the common Tutte/core/
centralizer upper-bound route during collaboration; the authoring
agent supplied the normalizer parity, overgroup analysis and explicit
type maps. The combined proof requires distinct review. It claims no
whole-source completion, new open-conjecture answer, Lean verification
or compression measurement.

\begin{thebibliography}{9}
\bibitem{source} F. Buekenhout, J. De Saedeleer and D. Leemans,
\emph{On the rank two geometries of the groups $\PSL(2,q)$: part II},
Ars Mathematica Contemporanea \textbf{6} (2013), 365--388,
Table 3 and p.~386.
\bibitem{tutte47} W. T. Tutte, \emph{A family of cubical graphs},
Proc. Cambridge Philos. Soc. \textbf{43} (1947), 459--474.
\bibitem{tutte59} W. T. Tutte, \emph{On the symmetry of cubic graphs},
Canadian J. Math. \textbf{11} (1959), 621--624.
\end{thebibliography}
\end{document}
